\begin{center}
    2. КАРТА РАСПРЕДЕЛЕНИЯ ПАМЯТИ
\end{center}

В \hyperref[tab:mem_alloc]{таблице~\ref*{tab:mem_alloc}} представлена карта распределения памяти.
В ней указаны исходные десятичные значения операндов, их шестнадцатеричные эквиваленты, 
а также базовые адреса загрузки данных и кода программы.

\begin{table}[!ht]
    \caption{Карта распределения памяти}
    \label{tab:mem_alloc}
    \centering
    \begin{tabular}{|c|c|c|c|}
        \hline
        Операнд & Десятичное значение & Шестнадцатеричный код & \multicolumn{1}{c|}{\begin{tabular}[c]{@{}c@{}}Адрес загрузки\\ (16)\end{tabular}} \\ \hline
        X\textsubscript{1}      & 18                 & 12                    & 00000008                                                                           \\ \hline
        X\textsubscript{2}      & --23               & E9                   & 00000012                                                                          \\ \hline
        X\textsubscript{3}      & 1225               & 4C9                   & 0000001C                                                                           \\ \hline
        X\textsubscript{4}      & --841              & FCB7                 & 00000026                                                                          \\ \hline
        X\textsubscript{5}      & 15000625           & E4E431                & 00000030                                                                           \\ \hline
        X\textsubscript{6}      & --707281           & FFF5352F             & 0000003A                                                                          \\ \hline
                &                    & Промежуточные ячейки  & 0000006С                                                                           \\ \hline
                &                    & Текст программы       & 2A8                                                                                \\ \hline
    \end{tabular}
\end{table}